% The Python chapter of the WFDB-x user's manual. It is the template of the chapters of
% the other languages, which follow the same sections in the same order. Each \apientry
% of the section "The high-level API" is checked by check_manual.py against the code
% (python/src/wfdbx): its signature listing must equal the name of the function followed
% by str(inspect.signature(...)), its params must name every parameter and no other, and
% every listing of style "example" is run (in a copy of tests/data).
\chapter{Python}
\label{chap:python}
\chapterlanguage{python}

The Python implementation of WFDB-x is the package \code{wfdbx}, published on PyPI as
\code{wfdb-x}. It is the reference implementation: the six other languages are
translations of it, and its results are those that they must reproduce. It is written in
Python alone, on top of NumPy, and never calls the C library.

It offers WFDB-x in three ways:
\begin{itemize}
\item \textbf{The high-level API} (\cref{sec:python-api}): functions that read and write
  whole records, headers, annotation files and header metadata as NumPy arrays and Python
  objects, in one call, with keyword arguments, and that raise exceptions on errors. They
  read signal files in large blocks, and are the fastest way to read a record. They are an
  addition of WFDB-x: the C package has no equivalent. They are the recommended interface
  for new code, analysis and notebooks.
\item \textbf{The WFDB library functions} under their C names (\code{isigopen},
  \code{getvec}, \code{getann}, \code{putann}, \ldots), with the behavior of the C library
  (\cref{sec:python-library}). They are documented once, for all the languages, in
  \cref{chap:library}.
\item \textbf{The command-line tools} (\code{rdsamp}, \code{rdann}, \code{wrsamp},
  \code{xform} and 29 others), with the options, output and exit status of the C tools,
  run through the \code{wfdbx} command (\cref{sec:python-tools}). They are documented in
  \cref{chap:tools}.
\end{itemize}
The high-level functions return exactly the samples and annotations that the library
functions (and the C library) give.

% ==========================================================================================
\section{Installation and requirements}
\label{sec:python-install}

\subsection{Requirements}

WFDB-x for Python needs Python 3.9 or later and NumPy 1.17 or later, its only dependency.
The package is pure Python: it has no compiled code and needs no C compiler, no WFDB
library and no libFLAC (the FLAC formats 508, 516 and 524 are read and written by
WFDB-x's own codec). It runs on Linux, macOS and Windows.

Two sets of optional dependencies are declared:
\begin{itemize}
\item \code{examples}: matplotlib and Jupyter, to run the example notebooks;
\item \code{test}: pytest, nbformat, nbclient and ipykernel, to run the tests
  (\cref{sec:python-testing}).
\end{itemize}

\subsection{Installing from PyPI}

\begin{lstlisting}[style=shell]
pip install wfdb-x                # the package and the wfdbx command
pip install "wfdb-x[examples]"    # also matplotlib and Jupyter, for the notebooks
\end{lstlisting}

The package is imported as \code{wfdbx} (\code{import wfdbx}); the installation also
creates the command \cmd{wfdbx}, which runs the command-line tools. The installed version
is \code{wfdbx.\_\_version\_\_}, or the output of \cmd{wfdbx -{}-version}:

\begin{lstlisting}[style=shell]
python -c "import wfdbx; print(wfdbx.__version__)"    # 1.2.2
wfdbx --version                                       # WFDB-x 1.2.2
\end{lstlisting}

The wheel and the source distribution include the licenses, \file{LICENSE} (the library,
LGPL-2.1-or-later) and \file{LICENSE-GPL} (the command-line tools, GPL-3.0-or-later), and
the file \file{NOTICE}, which says which parts are under which license and gives the
BSD-3-Clause notice of libFLAC, from which the FLAC decoder (\file{wfdbx/\_flac.py}) is
translated (\cref{sec:license}). The license expression of the package is
\code{LGPL-2.1-or-later AND GPL-3.0-or-later AND BSD-3-Clause}.

\subsection{Installing from the repository}

From the root of a copy of the repository, the package is in the directory
\file{python}:

\begin{lstlisting}[style=shell]
pip install ./python                  # install the package
pip install -e "./python[test]"       # or: editable, with the test dependencies
\end{lstlisting}

Without installation, the package can also be used from its source directory by adding
\file{python/src} to \code{PYTHONPATH}; the command-line tools are then run as
\cmd{python -m wfdbx.cli} (\cref{sec:python-tools}).

\subsection{The sample record and the example notebooks}
\label{sec:python-resources}

The package includes the sample record \code{100s}, the sample record of the WFDB
Software Package: the first minute of record 100 of the MIT-BIH Arrhythmia Database, two
ECG signals (\code{MLII} and \code{V5}) sampled at 360\,Hz, with the reference beat
annotations (annotator \code{atr}). Its files (\file{100s.hea}, \file{100s.dat},
\file{100s.atr}) are in the directory returned by \code{wfdbx.sample\_data\_dir()}, so
that the examples and the notebooks run without a network connection. Work that uses the
MIT-BIH Arrhythmia Database cites it as described on PhysioNet.

The package also includes seven Jupyter notebooks (\cref{tab:python-notebooks}), in the
directory returned by \code{wfdbx.examples\_dir()}. Each runs without a network
connection (\code{06\_physionet\_records} reads record 100 from PhysioNet when it can,
and the sample record otherwise). The function \code{copy\_examples}, or the command
\cmd{wfdbx -{}-examples DIR}, copies them to a directory where they can be opened and run:

\begin{lstlisting}[style=shell]
wfdbx --examples my-examples      # without DIR: ./wfdbx-examples
cd my-examples
jupyter lab
\end{lstlisting}

\begin{table}[htbp]
\centering
\caption{The example notebooks of the Python package}
\label{tab:python-notebooks}
\begin{tabularx}{\linewidth}{@{}lL@{}}
\toprule
\tabhead Notebook & \tabhead Contents \\
\midrule
\code{01\_getting\_started} & reading headers and signals; ADC and physical units \\
\code{02\_annotations\_and\_heart\_rate} & reading annotations; RR intervals and heart
  rate \\
\code{03\_writing\_records} & writing records and annotations \\
\code{04\_command\_line\_tools} & running the tools \\
\code{05\_c\_library\_api} & the WFDB library functions \\
\code{06\_physionet\_records} & reading a record from PhysioNet \\
\code{07\_legacy\_mode} & the original package's behavior, on request \\
\bottomrule
\end{tabularx}
\end{table}

% ==========================================================================================
\section{Quick start}
\label{sec:python-quickstart}

This example reads the header, ten seconds of the signals and the annotations of the
sample record \code{100s}, from the package's data directory (\code{input\_dir}), and
computes the mean heart rate. It then writes the first ten seconds as a new record,
\code{first10}, in the current directory, with two metadata entries in its header.

\begin{lstlisting}[language=Python,style=example]
import numpy as np
import wfdbx

sample = wfdbx.sample_data_dir()        # the directory of the sample record 100s

# The header only: no samples are read.
hdr = wfdbx.read_header("100s", input_dir=sample)
print(hdr.nsig, hdr.fs, hdr.nsamples)   # 2 360.0 21600
print([s.desc for s in hdr.siginfo])    # ['MLII', 'V5']

# Seconds 10 to 20 of both signals, in physical units (float64): one row per sample,
# one column per signal.
rec = wfdbx.read_record("100s", sampfrom="0:10", sampto="0:20", physical=True,
                        input_dir=sample)
print(rec.signals.shape, rec.signals.dtype)    # (3600, 2) float64
print(rec.signals[0])                          # [-0.39  -0.275]

# The beat annotations, and the mean heart rate.
ann = wfdbx.read_annotations("100s", "atr", input_dir=sample)
print(len(ann), ann.symbols[:3], ann.samples[:3])    # 75 ['+', 'N', 'N'] [ 18  77 370]
beats = ann.samples[[bool(wfdbx.isqrs(c)) for c in ann.codes]]
rr = np.diff(beats) / ann.fs                         # RR intervals, in seconds
print(len(beats), "beats,", round(60 / rr.mean(), 1), "beats per minute")   # 74 beats, 73.9 ...

# Write the first 10 seconds, in ADC units, as the record "first10" (first10.hea and
# first10.dat, in the current directory), with the signal specifications of 100s.
x = wfdbx.read_record("100s", sampto="0:10", input_dir=sample)
wfdbx.write_record("first10", x.signals, fs=x.fs, fmt=16, gains=x.gains,
                   baselines=x.baselines, adczeros=[s.adczero for s in x.siginfo],
                   adcres=[s.adcres for s in x.siginfo], units="mV", names=x.names,
                   metadata={"Source": "MIT-BIH record 100"})

# Header metadata: "# Key: value" lines of the header.
wfdbx.set_metadata("first10", "Age", "69")
print(wfdbx.read_metadata("first10"))   # [('Source', 'MIT-BIH record 100'), ('Age', '69')]
print(wfdbx.check_record("first10"))    # no problems found
\end{lstlisting}

The header \file{100s.hea} gives no units: the example writes \code{"mV"} explicitly,
since WFDB-x never assumes units, a gain, an ADC resolution or a sampling frequency that
a record does not give (\cref{chap:behavior}).

% ==========================================================================================
\section{The high-level API}
\label{sec:python-api}

The high-level API is made of the functions of this section, all available at the top
level of the package (\code{wfdbx.read\_record}, \ldots), and of the classes of their
results (\cref{sec:python-results}). They are defined in the modules
\code{wfdbx.highlevel}, \code{wfdbx.metadata}, \code{wfdbx.check} and
\code{wfdbx.resources}; \code{help()} gives the documentation of each.

\subsection{Conventions}
\label{sec:python-conventions}

\paragraph{Record names and directories.} A record is named as in the other parts of
WFDB-x (\cref{chap:records,sec:record-paths}), without the \file{.hea} extension (a
trailing \code{".hea"} is removed, by the readers and the writers). A record name is one
of:
\begin{itemize}
\item a full path, absolute (\code{"/data/mitdb/100"}) or in the home directory
  (\code{"\textasciitilde/Documents/100s"}), used as given;
\item a relative path (\code{"100s"}, \code{"data/100s"}), read from the input directory
  if one is set, otherwise from the current directory, and written to the output
  directory if one is set, otherwise to the current directory (missing directories are
  created);
\item a URL (\code{"https://physionet.org/files/mitdb/1.0.0/100"}), read only if network
  access is on; a URL cannot be written.
\end{itemize}
The input and output directories are set with \code{wfdbx\_set\_input\_dir} and
\code{wfdbx\_set\_output\_dir} (\cref{sec:python-settings}) or with the environment
variables \code{WFDBX\_INPUT\_DIR} and \code{WFDBX\_OUTPUT\_DIR}. The functions that read
take the optional keyword arguments \code{input\_dir} and \code{network}, and those that
write take \code{dir}: they apply to that one call, and the previous settings are
restored when it returns, also after an error. \code{input\_dir} applies to relative
names only, not to full paths or URLs. A writer given a name with a path creates the
files in that directory, and the header names the record and its signal file without the
path (\code{rec}, \code{rec.dat}), so that the record opens under any path that leads to
it; in legacy mode, as in the original, the header gives the name with its path. With
\code{dir}, the name must be a plain name (\code{"rec"}): the files are created in
\code{dir}, and the record opens from that directory under its plain name. The signal
files that a header names are read from the header's directory.

\paragraph{Times.} The arguments \code{sampfrom} and \code{sampto} are sample numbers
(whole numbers) or WFDB time strings, as \code{strtim} reads them: \code{"0:10"} (10
seconds), \code{"1:02:03.5"}, \code{"s3600"} (sample 3600), \code{"e"} (the end of the
record), or a time of day in brackets, \code{"[14:00:00]"}. A time in seconds needs the
sampling frequency; when the record does not give it, such a time raises
\code{ValueError}.

\paragraph{Per-signal and per-annotation arguments.} The arguments that describe each
signal (\code{fmt}, \code{gains}, \code{units}, \ldots) or each annotation
(\code{subtypes}, \code{chans}, \ldots) take one value for all, or a list (or tuple, or
NumPy array) with exactly one value per signal or per annotation.

\paragraph{Signal parameters.} No sampling frequency, gain, ADC resolution or units is
ever assumed. Where one is needed and the record does not give it, it must be given as an
argument; otherwise the function raises \code{ValueError}. A record whose header gives no
sampling frequency has \code{fs = 0}.

\paragraph{Legacy mode.} With \code{legacy=True}, a function behaves as in the original
WFDB package, with its assumed signal parameters and its search for files
(\cref{chap:behavior}); the library's legacy mode is on for the duration of the call
only. With \code{legacy=False} (the default), the current setting of
\code{wfdbx\_set\_legacy} applies, which is off unless it was turned on.

\paragraph{State.} The functions that read and write records and annotation files close
anything opened with the WFDB library functions (they call \code{wfdbquit()}) before
they start, and close their own files before they return, whatever the outcome.
\code{read\_record} sets the getvec mode for its own reading only. \code{read\_record},
\code{read\_header}, \code{read\_annotations}, \code{write\_record},
\code{write\_annotations} and \code{check\_record} leave the time resolution of output
annotators (\code{getafreq()}) and the getvec mode (\code{getgvmode()}) as they were, also
after reading an EDF record, whose opening sets both (as in C; the library functions keep
that behavior). The metadata functions open no signals and leave the library's state as
it is.

\paragraph{The examples} of this section read the test records of WFDB-x (among them
\code{100s}) from the current directory, and write records of their own there. To run
them elsewhere, copy the sample record to the current directory first, for example with
\code{shutil.copytree(wfdbx.sample\_data\_dir(), ".", dirs\_exist\_ok=True)}.

% ------------------------------------------------------------------------------------------
\subsection{Reading records}
\label{sec:python-reading}

\apientry{read\_record}{Read the signals of a record}
\begin{lstlisting}[language=Python,style=signature]
read_record(record, channels=None, sampfrom=0, sampto=None, physical=False, highres=False, pad=False, *, input_dir=None, network=None, legacy=False, checksum_warning=True)
\end{lstlisting}
Reads the signals of a record, or part of it, into a NumPy array with one row per sample
and one column per signal, together with the description of the record. The samples are
those that \code{getvec} returns (\cref{chap:library}), read in large blocks. Invalid
samples are \code{wfdbx.WFDB\_INVALID\_SAMPLE} ($-32768$) in ADC units, and NaN in
physical units. When the read reaches the end of the record, the samples are compared
with the checksums of the header (\cref{sec:python-warnings}). The result is allocated
for at most $2^{26}$ samples in advance, whatever length the header gives, and grows as
the samples are read.
\begin{params}
\item[record] Record name (a string), without the \file{.hea} extension (a trailing
  \code{".hea"} is also accepted): a full path, absolute (\code{"/data/mitdb/100"}) or in
  the home directory (\code{"\textasciitilde/Documents/100s"}), used as given; a relative
  path (\code{"100s"}, \code{"data/100s"}), read from the input directory if one is set,
  otherwise from the current directory; or a URL
  (\code{"https://physionet.org/files/mitdb/1.0.0/100"}), read only if network access is
  on.
\item[channels] The signals to read, in the order in which they are returned: a list of
  signal numbers (\code{int}, from 0) or signal descriptions (\code{str}, as in
  \code{Record.names}, such as \code{"MLII"}), which may be mixed. Default \code{None}:
  all the signals, in the order of the header.
\item[sampfrom] The first sample to read: a sample number (a whole number $\ge 0$) or a
  time string. Default \code{0}: the beginning of the record.
\item[sampto] The sample after the last one to read: a sample number (a whole number) or
  a time string. A value at or before \code{sampfrom} reads no samples. Default
  \code{None}: the end of the record.
\item[physical] \code{False} (the default) to return ADC units, as \code{int32};
  \code{True} to return physical units, as \code{float64}, computed as
  $(\mathit{adc} - \mathit{baseline}) / \mathit{gain}$. Every signal read must then have a
  gain.
\item[highres] \code{False} (the default) to read a multi-frequency record at its frame
  frequency, with the samples of the faster signals averaged to one per frame;
  \code{True} to read every sample, at the highest sampling frequency of the record (the
  samples of the slower signals are repeated). \code{Record.fs} is then that frequency.
\item[pad] \code{True} to replace each invalid sample by the previous valid sample of the
  same signal. Default \code{False}.
\item[input\_dir] The input directory for this call only (optional): a local directory, a
  \code{\textasciitilde} path or the URL of a directory, under which relative record names
  are read. It does not apply to full paths or URLs. Default \code{None}: the current
  setting (\code{wfdbx\_set\_input\_dir}).
\item[network] \code{True} to allow reading URLs for this call only, \code{False} to forbid
  it. Default \code{None}: the current setting (\code{wfdbx\_set\_network}).
\item[legacy] \code{True} to behave as the original WFDB package (legacy mode, for this
  call only). Default \code{False}.
\item[checksum\_warning] \code{True} (the default) to issue a \code{RuntimeWarning} when
  the samples do not match the checksums of the header; \code{False} for no warning. The
  samples are returned either way, and \code{checksum\_ok} of the result is \code{False}.
\end{params}
\returns{A \code{Record} (\cref{tab:python-record}).}
\raises{\code{IOError} (that is, \code{OSError}) if the header or a signal file cannot be
read, if the record has no signals, if \code{sampfrom} is beyond the end of the record,
or if a signal file is shorter than its header says or cannot be read (``can't read
signal data of record REC (getvec status -3)''), even if the checksums do not match
either: no incomplete frame is ever returned; \code{ValueError} if \code{sampfrom} is
negative, if \code{sampfrom} or \code{sampto} is not a whole number (a fraction, NaN or an infinite
value), if a channel is not in the record, if a time string needs the sampling frequency
and the header gives none, or if \code{physical=True} and a signal read has no gain.
Warns with a \code{RuntimeWarning} on a checksum error (see \code{checksum\_warning}).}
\begin{lstlisting}[language=Python,style=example]
import wfdbx

# Seconds 10 to 20 of lead MLII, in physical units.
rec = wfdbx.read_record("100s", channels=["MLII"], sampfrom="0:10", sampto="0:20",
                        physical=True)
print(rec.signals.shape, rec.signals.dtype)   # (3600, 1) float64
print(rec.names, rec.fs, rec.sampfrom)        # ['MLII'] 360.0 3600

# The first 3 samples of both signals, in ADC units, V5 first.
rec = wfdbx.read_record("100s", channels=[1, 0], sampto=3)
print(rec.signals.tolist())     # [[1011, 995], [1011, 995], [1011, 995]]
print(rec.names, rec.gains, rec.baselines, rec.units)
# ['V5', 'MLII'] [200.0, 200.0] [1024, 1024] [None, None]

# The whole record: the checksums are verified.
rec = wfdbx.read_record("100s")
print(rec.signals.shape, rec.checksum_ok)     # (21600, 2) True
\end{lstlisting}
A multi-frequency record (here the test record \code{mfreq}, whose signals have 1, 2
and 4 samples per frame at a frame frequency of 100\,Hz) is read at its frame frequency,
or at its highest sampling frequency:
\begin{lstlisting}[language=Python,style=example]
import wfdbx

rec = wfdbx.read_record("mfreq")
print(rec.fs, rec.signals.shape)       # 100.0 (600, 3)
rec = wfdbx.read_record("mfreq", highres=True)
print(rec.fs, rec.signals.shape)       # 400.0 (2400, 3)
print([s.spf for s in rec.siginfo])    # [1, 2, 4]
\end{lstlisting}

\apientry{read\_header}{Read the header of a record}
\begin{lstlisting}[language=Python,style=signature]
read_header(record, *, input_dir=None, network=None, legacy=False)
\end{lstlisting}
Reads the header of a record, without opening its signal files. The result describes the
record and each of its signals, and holds the info strings and the metadata entries of the
header. For a multi-segment record, the master header is read, and \code{nsamples} is the
length of the whole record.
\begin{params}
\item[record] Record name (a string), as for \code{read\_record}: a full path (absolute or
  \code{\textasciitilde}), used as given; a relative path, read from the input directory
  if one is set, otherwise from the current directory; or a URL, read only if network
  access is on.
\item[input\_dir] The input directory for this call only (optional; as for
  \code{read\_record}), under which relative record names are read; it does not apply to
  full paths or URLs. Default \code{None}: the current setting.
\item[network] \code{True} to allow reading URLs for this call only, \code{False} to forbid
  it. Default \code{None}: the current setting.
\item[legacy] \code{True} to read the header as the original WFDB package does. Default
  \code{False}.
\end{params}
\returns{A \code{Header} (\cref{tab:python-header}).}
\raises{\code{IOError} if the header cannot be found or read.}
\begin{lstlisting}[language=Python,style=example]
import wfdbx

hdr = wfdbx.read_header("100s")
print(hdr.name, hdr.fs, hdr.nsig, hdr.nsamples)    # 100s 360.0 2 21600
for s in hdr.siginfo:
    print(s.desc, s.fname, s.fmt, s.gain, s.adcres, s.units)
# MLII 100s.dat 212 200.0 11 None
# V5 100s.dat 212 200.0 11 None
print(hdr.comments[0])          # " 69 M 1085 1629 x1": the first info string
print(hdr.base_time)            # None: the header gives no start time
\end{lstlisting}

\apientry{read\_annotations}{Read an annotation file}
\begin{lstlisting}[language=Python,style=signature]
read_annotations(record, annotator, sampfrom=0, sampto=None, fs=None, *, input_dir=None, network=None, legacy=False)
\end{lstlisting}
Reads the annotations of an annotation file, in file order, into NumPy arrays and lists.
The sampling frequency, needed to convert times in seconds, is that of the record's
header, or, for a record without a header or whose header gives none, the argument
\code{fs}.
\begin{params}
\item[record] Record name (a string), as for \code{read\_record}: a full path (absolute or
  \code{\textasciitilde}), used as given; a relative path, read from the input directory
  if one is set, otherwise from the current directory; or a URL, read only if network
  access is on. Its header, if there is one, gives the sampling frequency.
\item[annotator] Annotator name: the suffix of the annotation file. \code{"atr"} reads the
  file \file{RECORD.atr}.
\item[sampfrom] The first time to read: a sample number (a whole number $\ge 0$) or a
  time string. Default \code{0}: the beginning of the file.
\item[sampto] The time at which to stop: annotations at or after it are not read. A
  sample number (a whole number) or a time string. Default \code{None}: the end of the
  file.
\item[fs] The sampling frequency, in Hz (a number $> 0$), to use when the record has no
  header or its header gives none; it is ignored otherwise. Default \code{None}: none.
\item[input\_dir] The input directory for this call only (optional; as for
  \code{read\_record}), under which relative record names are read; it does not apply to
  full paths or URLs. Default \code{None}: the current setting.
\item[network] \code{True} to allow reading URLs for this call only, \code{False} to forbid
  it. Default \code{None}: the current setting.
\item[legacy] \code{True} to behave as the original WFDB package (which assumes 250\,Hz
  for a record without a frequency). Default \code{False}.
\end{params}
\returns{An \code{Annotations} object (\cref{tab:python-annotations}). Its \code{fs} is
the frequency in which the times are counted: the one written in the annotation file, if
there is one, or else the sampling frequency; it is 0 if none of them is known.}
\raises{\code{IOError} if the annotation file cannot be opened or read, or if it ends
before its end-of-file marker (``can't read annotation file ANN of record REC (getann
status -3)''); \code{ValueError} if \code{fs} is given and is not $> 0$, if
\code{sampfrom} is negative (``sampfrom must not be negative''), if \code{sampfrom} or
\code{sampto} is not a whole number (a fraction, NaN or an infinite value: ``sampfrom must
be a whole number''), or if a time string needs the sampling frequency and none is known.
The arguments are checked before the annotation file is opened (a time string is converted
with the sampling frequency of the record's header, which is read for it), and the
annotation file is closed whatever happens.}
\begin{lstlisting}[language=Python,style=example]
import wfdbx

ann = wfdbx.read_annotations("100s", "atr", sampfrom="0:10", sampto="0:30")
print(len(ann), ann.fs)                   # 24 360.0
print(ann.samples[:3], ann.symbols[:3])   # [3862 4170 4466] ['N', 'N', 'N']

# The whole file: the first annotation is a rhythm change, with an aux string.
ann = wfdbx.read_annotations("100s", "atr")
print(ann.symbols[0], ann.codes[0] == wfdbx.RHYTHM, ann.aux[0])   # + True (N
print(sorted(set(ann.symbols)))           # ['+', 'A', 'N']
\end{lstlisting}
For a record without a header, \code{fs} gives the frequency in which times in seconds
are converted (the library reports on the standard error that it cannot open the
header):
\begin{lstlisting}[language=Python,style=example]
import wfdbx

wfdbx.write_annotations("beats", "qrs", [100, 460, 820], symbols="N")   # no header
ann = wfdbx.read_annotations("beats", "qrs", sampto="0:02", fs=360)
print(ann.samples, ann.fs)                # [100 460] 360.0
\end{lstlisting}

% ------------------------------------------------------------------------------------------
\subsection{Writing records}
\label{sec:python-writing}

\apientry{write\_record}{Write a header and a signal file}
\begin{lstlisting}[language=Python,style=signature]
write_record(record, signals, fs, fmt=16, gains=None, baselines=None, adczeros=0, adcres=None, units=None, names=None, base_time=None, comments=None, physical=False, samps_per_frame=None, *, metadata=None, dir=None, legacy=False)
\end{lstlisting}
Writes signals as a new record: the header \file{RECORD.hea} and one signal file,
\file{RECORD.dat}, which holds all the signals (so they share one format). Existing files
of these names are replaced: both files are written to temporary files in the
destination directory (\file{.NAME.XXXXXXXX.tmp}, readable and writable by the owner
only), which replace the existing files, with their permissions, only when both are
complete; after an error, the temporary files are removed and the existing files are left
as they were. For a symbolic link, the file it leads to is replaced; a destination that
is not a regular file (a device) is written directly. The sampling frequency, the gain and
ADC resolution of each signal, and the units of each calibrated signal must be given:
there are no defaults. The per-signal arguments (\code{fmt}, \code{gains}, \code{baselines},
\code{adczeros}, \code{adcres}, \code{units}, \code{names}) take one value for all the
signals, or a list with one value per signal. All the arguments are checked before any
file is created or changed.
\begin{params}
\item[record] Record name (a non-empty string), without the \file{.hea} extension (a
  trailing \code{".hea"} is removed, as by the readers): a plain name (\code{"rec"}) or a relative path (\code{"data/rec"}), created in the output
  directory if one is set, otherwise in the current directory, missing directories
  included; or a full path, absolute or in the home directory
  (\code{"\textasciitilde/Documents/rec"}), used as given. With a path, the files are
  created in that directory, and the header names the record and its signal file without
  the path (\code{rec}, \code{rec.dat}), so that the record opens under any path that
  leads to it (in legacy mode, as in the original, the header gives the name with its
  path). A URL cannot be written.
\item[signals] The samples: an array (or anything \code{numpy.asarray} accepts) with one
  row per sample and one column per signal (2-D); a 1-D array is one signal. In ADC units
  (the default), the values are integers, or floats with whole values, within the range of
  a 32-bit integer; with \code{physical=True}, they are in physical units.
\item[fs] The sampling frequency, in Hz: a finite number $> 0$. Required.
\item[fmt] The storage format, such as \code{16}, \code{212}, \code{24} or \code{516}
  (FLAC); one value, or one per signal, all equal (\cref{chap:formats}). Default
  \code{16}.
\item[gains] The gain of each signal, in ADC units per physical unit: a finite number;
  \code{0} marks an uncalibrated signal, which is allowed only for samples given in ADC
  units. Negative gains are allowed. Required (default \code{None}: not given).
\item[baselines] The ADC value that corresponds to 0 physical units (a whole number
  within the range of a 32-bit integer). Default \code{None}: the ADC zero of the signal.
\item[adczeros] The ADC value at the middle of the ADC range (a whole number within the
  range of a 32-bit integer). Default \code{0}.
\item[adcres] The ADC resolution, in bits (a whole number $\ge 1$). Required (default
  \code{None}: not given).
\item[units] The physical units, such as \code{"mV"} (Latin-1 text, one byte per
  character, as all text in \code{wfdbx}). Required for each signal whose gain is not 0.
  Default \code{None}: not given.
\item[names] The signal descriptions, such as \code{"MLII"} (Latin-1 text). Default
  \code{None}: no descriptions.
\item[base\_time] The start time, with an optional date:
  \code{"hh:mm:ss"} or \code{"hh:mm:ss dd/mm/yyyy"}, with optional decimals of the
  seconds. Default \code{None}: no start time.
\item[comments] A list of info strings (Latin-1 text), written to the header as comment lines
  (\code{"\#"} followed by the string). Default \code{None}: none.
\item[physical] \code{True} if \code{signals} are in physical units. They are converted
  to ADC units as $\mathrm{round}(x \cdot \mathit{gain}) + \mathit{baseline}$ (halves
  rounded away from zero); NaN becomes an invalid sample, and infinite values and values
  beyond the range of 32-bit ADC units are errors. Default \code{False}: ADC units.
\item[samps\_per\_frame] For a multi-frequency record, the number of samples per frame of
  each signal (integers $\ge 1$). \code{signals} then has one row per frame and
  \code{sum(samps\_per\_frame)} columns: the samples of signal 0, then those of signal 1,
  and so on. \code{fs} is then the frame frequency. Default \code{None}: one sample per
  frame for every signal.
\item[metadata] Metadata entries (\cref{chap:metadata}) to write to the header after the
  comments, as \code{\# Key: value}: a mapping from key to value (strings), or a sequence
  of (key, value) pairs. Keys and values follow the rules of \code{set\_metadata}; no two
  keys may differ only in letter case, nor be the key of an entry in \code{comments}.
  Default \code{None}: none.
\item[dir] The directory (which must exist) in which to create the files, for this call
  only. \code{record} must then be a plain name, without a directory, and the record opens
  from that directory under its plain name. Default \code{None}: the output directory.
\item[legacy] \code{True} to behave as the original WFDB package: the gain defaults to
  200, and the ADC resolution to one based on the format. Default \code{False}.
\end{params}
\returns{\code{None}.}
\raises{\code{ValueError} if \code{record} is empty; if \code{signals} is not a 1-D or 2-D
array (``signals must be a 1-D or 2-D array''); if it has no signals
(``at least one signal is required''); if \code{fs}, a gain, an ADC resolution or the
units of a calibrated signal is missing or invalid (\code{fs} and the gains must be
finite, and \code{fs} $> 0$); if a baseline, an ADC zero or an ADC resolution is not a
whole number within the range of a 32-bit integer (``baselines must be whole numbers
within the range of a 32-bit integer''); if the record name, a signal name, the units or
a comment has a character that is not Latin-1 (``units 'XX' has a character that is not
Latin-1 (one byte per character)''); if \code{base\_time} cannot be parsed; if a list of
per-signal values does not have one value per signal; if a signal in physical units has a
gain of 0; if a sample is not valid (see \code{signals} and \code{physical}); if the
signals do not all have the same format; if a \code{samps\_per\_frame} value is not an
integer $\ge 1$, or the number of columns of \code{signals} is not their sum; if an entry
of \code{metadata} is not valid; or if \code{record} contains a directory and \code{dir}
is given. \code{TypeError} if \code{metadata} is neither a mapping nor a sequence of
pairs, or a key or value is not a string. \code{FileNotFoundError} if \code{dir} does not
exist. \code{IOError} if the signal file or the header cannot be created, or cannot be
written completely, for example on a full disk (``write error on record REC'').}
\begin{lstlisting}[language=Python,style=example]
import numpy as np
import wfdbx

# Two signals of random numbers, in millivolts, at 360 Hz: the record "rand"
# (rand.hea and rand.dat).
x = np.random.default_rng(1).standard_normal((3600, 2))
wfdbx.write_record("rand", x, fs=360, fmt=16, physical=True, gains=1000, adcres=16,
                   units="mV", names=["A", "B"], base_time="14:00:03 21/06/2024",
                   comments=["random numbers"], metadata=[("Generator", "numpy")])
rec = wfdbx.read_record("rand", physical=True)
print(rec.signals.shape, rec.base_time)       # (3600, 2) [14:00:03 21/06/2024]
print(np.abs(rec.signals - x).max() <= 0.0005)        # True: rounded to 1/1000 mV
print(wfdbx.read_header("rand").comments)     # ['random numbers', ' Generator: numpy']

# An uncalibrated signal (gain 0, no units), in ADC units, in the directory "out".
import os
os.makedirs("out", exist_ok=True)
wfdbx.write_record("counts", np.arange(1000), fs=100, gains=0, adcres=12, dir="out")
print(wfdbx.read_record("out/counts").signals[:3, 0])    # [0 1 2]
\end{lstlisting}
A multi-frequency record: signal \code{fast} has 2 samples per frame, signal
\code{slow} one; each row of \code{signals} is one frame, with the two samples of
\code{fast} and then the sample of \code{slow}:
\begin{lstlisting}[language=Python,style=example]
import numpy as np
import wfdbx

frames = 100                                    # 1 second at 100 frames per second
fast = np.arange(2 * frames).reshape(frames, 2)     # 200 samples per second
slow = np.full((frames, 1), 500)                    # 100 samples per second
wfdbx.write_record("mf", np.hstack([fast, slow]), fs=100, gains=200, adcres=12,
                   units="mV", names=["fast", "slow"], samps_per_frame=[2, 1])
rec = wfdbx.read_record("mf", highres=True)
print(rec.fs, rec.signals[:3].tolist())    # 200.0 [[0, 500], [1, 500], [2, 500]]
\end{lstlisting}

\apientry{write\_annotations}{Write an annotation file}
\begin{lstlisting}[language=Python,style=signature]
write_annotations(record, annotator, samples, symbols=None, codes=None, subtypes=0, chans=0, nums=0, aux=None, fs=None, *, sort=None, dir=None, legacy=False)
\end{lstlisting}
Writes an annotation file, \file{RECORD.ANNOTATOR}, in the MIT format. An existing file of
that name is replaced: the new file is written to a temporary file in the same directory
(readable and writable by the owner only), which replaces it, with its permissions, only
when it is complete (and, if needed, sorted); after an error, the temporary file is
removed and the existing file is left as it was. The record need not have a header. The
annotation types are given as \code{symbols} or as \code{codes} (\cref{chap:formats}). Each per-annotation
argument is one value for all the annotations, or a list with exactly one value per
annotation. If the annotations are not in time order, the file is sorted when it is
written, unless \code{sort=False}. All the arguments are checked before the file is
created or changed.
\begin{params}
\item[record] Record name (a string), without the \file{.hea} extension (a trailing
  \code{".hea"} is also accepted): a plain name or a relative path (\code{"data/100s"}),
  created in the output directory if one is set, otherwise in the current directory; or a
  full path, absolute or in the home directory (\code{"\textasciitilde/Documents/100s"}),
  used as given. The file is \file{RECORD.ANNOTATOR} at that place
  (\file{data/100s.qrs}). A URL cannot be written.
\item[annotator] Annotator name: the suffix of the file, made of letters, digits,
  \code{\_}, \code{-} and \code{\textasciitilde}.
\item[samples] The annotation times, in samples: a 1-D array or sequence of integers, one
  per annotation, from $-2^{63}+1$ to $2^{63}-2$ (the smallest and largest 64-bit
  integers, \code{WFDB\_TIME\_MIN} and \code{WFDB\_TIME\_MAX}, cannot be written).
\item[symbols] The mnemonics of the annotation types, such as \code{"N"} or \code{"V"}:
  one per annotation, or one string for all. An unknown mnemonic becomes code 0
  (\code{NOTQRS}). Default \code{None}: the types are given by \code{codes}.
\item[codes] The annotation codes (\code{wfdbx.NORMAL}, \code{wfdbx.PVC}, \ldots), from
  0 to 58 (59 to 63 are the control codes of the file format, \cref{sec:formats-mit-annotations}),
  instead of \code{symbols}: one per annotation, or one value for all. If both are given
  (as \code{read\_annotations} returns them), each symbol must be that of its code (an
  empty symbol is not checked). Default \code{None}.
\item[subtypes] The subtype of each annotation, stored as a signed byte (other values wrap
  around to $-128 \ldots 127$, as in C). Default \code{0}.
\item[chans] The channel number of each annotation, stored as an unsigned byte (other
  values wrap around to $0 \ldots 255$). Default \code{0}.
\item[nums] The annotator number of each annotation, stored as a signed byte. Default
  \code{0}.
\item[aux] The auxiliary text of each annotation, with \code{None} where there is none:
  a list, or one text for all. Texts longer than 255 bytes are cut. Default \code{None}:
  none.
\item[fs] If given (a finite number $> 0$), the frequency in which the times are
  counted, written to the file; it applies to this file only. Default \code{None}: none is
  written.
\item[sort] \code{True} to always sort the file, \code{False} to keep the given order
  (without the library's advice to run \code{sortann}), for this call only. Default
  \code{None}: the environment variable \code{WFDBANNSORT} decides (sorting is on unless
  it is \code{0}).
\item[dir] The directory (which must exist) in which to create the file, for this call
  only. \code{record} must then be a plain name, without a directory: the file is
  \file{DIR/RECORD.ANNOTATOR}. Default \code{None}: the output directory.
\item[legacy] \code{True} to behave as the original WFDB package. Default \code{False}.
\end{params}
\returns{\code{None}.}
\raises{\code{ValueError} if neither \code{symbols} nor \code{codes} is given for one
annotation or more (\code{symbols or codes must be given}; with no annotations, an empty
file is written); if both are
given and disagree (``symbols and codes disagree at annotation K''); if \code{samples}
(or a sequence of codes, subtypes, channels or annotator numbers) is not 1-D (``samples
must be a 1-D array''); if a list of per-annotation values does not have one value per
annotation; if a time is not an integer
within the range of a 64-bit integer (NaN, an infinite value, a fraction or a huge
value), or is $-2^{63}$ or $2^{63}-1$ (``annotation time T cannot be written''); if a
code is not from 0 to 58 (``annotation code 61 is reserved (59 to 63 are the control
codes of the annotation file format)'', ``annotation code 64 is out of range (0 to
58)''); if a sample, code, subtype, channel or annotator number is NaN or \code{None}; if
an aux text has a character that is not Latin-1; if \code{fs} is not finite and $> 0$; or if \code{record} contains a directory and
\code{dir} is given. \code{FileNotFoundError} if \code{dir} does not exist.
\code{IOError} if the annotation file cannot be created, or cannot be written completely,
or the library (\code{putann}) reports an error while writing it (``write error on
annotator ANN for record REC'').}
\begin{lstlisting}[language=Python,style=example]
import wfdbx

# Three beats and a rhythm change, by symbol, with an aux string for the rhythm.
wfdbx.write_annotations("100s", "test", [18, 77, 370, 662],
                        symbols=["+", "N", "V", "N"], aux=["(N", None, None, None])
ann = wfdbx.read_annotations("100s", "test")
print(ann.symbols, ann.aux)     # ['+', 'N', 'V', 'N'] ['(N', None, None, None]

# A copy of the reference annotations, as read_annotations returns them.
a = wfdbx.read_annotations("100s", "atr")
wfdbx.write_annotations("100s", "copy", a.samples, symbols=a.symbols, codes=a.codes,
                        subtypes=a.subtypes, chans=a.chans, nums=a.nums, aux=a.aux)
b = wfdbx.read_annotations("100s", "copy")
print((b.samples == a.samples).all(), b.symbols == a.symbols, b.aux == a.aux)
# True True True

# Annotations out of time order are sorted when the file is written (the library
# reports it on the standard error).
wfdbx.write_annotations("100s", "unsorted", [500, 100], codes=wfdbx.NORMAL)
print(wfdbx.read_annotations("100s", "unsorted").samples)      # [100 500]

# A reserved code is refused before anything is written: the file is unchanged.
try:
    wfdbx.write_annotations("100s", "unsorted", [100, 200], codes=[wfdbx.NORMAL, 61])
except ValueError as e:
    print("ValueError:", e)   # ValueError: annotation code 61 is reserved (59 to 63 ...)
print(wfdbx.read_annotations("100s", "unsorted").samples)      # [100 500]
\end{lstlisting}

% ------------------------------------------------------------------------------------------
\subsection{Checking records}
\label{sec:python-checking}

\apientry{check\_record}{Check a record and its annotation files}
\begin{lstlisting}[language=Python,style=signature]
check_record(record, annotators=(), *, input_dir=None, network=None)
\end{lstlisting}
Checks the header of a record, its signal specifications, every sample of its signal
files, each segment of a multi-segment record, and the annotation files named, and
reports each problem found. A problem in the record never raises an exception: it is
reported, with a severity (\code{"error"}: the record is damaged or does not follow the
WFDB format, and programs will read wrong data or fail; \code{"warning"}: the record can
be read, but something is missing or unusual, and may give wrong results; \code{"info"}:
a fact worth knowing, such as invalid samples). The problems, their codes and their
messages are the same in the seven languages (\cref{tab:python-issue-codes}). The checks
run in normal mode, even if legacy mode is on, and the library prints no messages during
the call. They are similar to those of Benjamin Moody's \code{wfdbcheck}.
\begin{params}
\item[record] Record name (a string), as for \code{read\_record}: a full path (absolute or
  \code{\textasciitilde}), used as given; a relative path, read from the input directory
  if one is set, otherwise from the current directory; or a URL, read only if network
  access is on.
\item[annotators] The annotator names whose files are checked, such as \code{["atr"]}; one
  name may be given as a string. Default \code{()}: none.
\item[input\_dir] The input directory for this call only (optional; as for
  \code{read\_record}), under which relative record names are read; it does not apply to
  full paths or URLs. Default \code{None}: the current setting.
\item[network] \code{True} to allow reading URLs for this call only, \code{False} to forbid
  it. Default \code{None}: the current setting.
\end{params}
\returns{A \code{CheckReport} (\cref{tab:python-checkreport}): \code{report.ok} is
\code{True} if no problem is an error, \code{report.issues} lists the problems
(\code{Issue} objects, \cref{tab:python-issue}), and \code{print(report)} prints one line
per problem, or \code{no problems found}.}
\raises{\code{TypeError} if \code{record} or an annotator name is not a string;
\code{ValueError} if \code{record} is empty. A record that cannot be read is reported
(code \code{header\_unreadable}), not raised.}
\begin{lstlisting}[language=Python,style=example]
import numpy as np
import wfdbx

report = wfdbx.check_record("100s", ["atr"])
print(report.ok)          # True: no errors
print(report)
# warning: signal 0: the units are not specified
# warning: signal 1: the units are not specified

# A damaged record: one sample of a 12-bit signal is overwritten with 2048.
wfdbx.write_record("damaged", np.arange(100), fs=100, gains=100, adcres=12, units="mV",
                   names="ramp")
with open("damaged.dat", "r+b") as f:
    f.seek(20)                    # sample 10 (2 bytes per sample in format 16)
    f.write(b"\x00\x08")          # 0x0800 = 2048, little-endian
report = wfdbx.check_record("damaged")
print(report.ok)          # False
for issue in report.issues:
    print(issue.severity, issue.code, issue.where)
# error sample_range signal 0
# error checksum signal 0
print(report.lines()[0])
# error: signal 0: 1 sample is outside the ADC range (-2048 to 2047): sample 10 (value 2048)
\end{lstlisting}

% ------------------------------------------------------------------------------------------
\subsection{Header metadata}
\label{sec:python-metadata}

These functions read and change the metadata of a header: its info strings of the form
\code{\# Key: value}, such as \code{\# Age: 63} (\cref{chap:metadata}). In short: keys are
compared without regard to letter case; a key must not be empty and must not contain
\code{:}, \code{\#}, a line break or a NUL character, and a value must not contain a line
break or a NUL character; spaces at either end of keys and values are removed; text is
one byte per character (Latin-1), and \code{bytes} are also accepted, as given. A change
rewrites the header in one step, through a temporary file, keeping every other byte of
it; all the arguments are checked first, so an error leaves the header unchanged. Only
the \file{.hea} file is read (not a \file{.info} file); for a multi-segment record, it is
the master header, unless a segment is named. Metadata often holds information about
patients: WFDB-x stores it as plain text, exactly as given, and does not de-identify it.

\begin{phibox}
A header can hold identifying data about a patient (a name, a date of birth, a medical
record number). The metadata functions read and write such data as given. Records that
are shared must be de-identified first.
\end{phibox}

\apientry{read\_metadata}{Read the metadata of a record}
\begin{lstlisting}[language=Python,style=signature]
read_metadata(record, *, input_dir=None, network=None)
\end{lstlisting}
Reads the metadata entries of the header of a record, in file order. An info string is an
entry if it contains a colon: the key is the text before the first colon, the value the
rest of the line, both without spaces at either end. Info strings without a colon are free
comments, which are not returned.
\begin{params}
\item[record] Record name (a string), without the \file{.hea} extension (a trailing
  \code{".hea"} is also accepted): a full path, absolute (\code{"/data/mitdb/100"}) or in
  the home directory (\code{"\textasciitilde/Documents/100s"}), used as given; a relative
  path (\code{"100s"}, \code{"data/100s"}), read from the input directory if one is set,
  otherwise from the current directory; or a URL
  (\code{"https://physionet.org/files/mitdb/1.0.0/100"}), read only if network access is
  on.
\item[input\_dir] The input directory for this call only (optional): a local directory, a
  \code{\textasciitilde} path or a URL, under which relative record names are read. It
  does not apply to full paths or URLs. Default \code{None}: the current setting.
\item[network] \code{True} to allow reading URLs for this call only, \code{False} to
  forbid it. Default \code{None}: the current setting.
\end{params}
\returns{A list of (key, value) pairs of strings, in file order. A key that appears more
than once (ignoring letter case) is returned once, with its first value, and a
\code{UserWarning} is issued (\cref{sec:python-warnings}).}
\raises{\code{IOError} if the header cannot be found or read; \code{ValueError} if
\code{record} is empty or \code{"-"} (the standard input).}
\begin{lstlisting}[language=Python,style=example]
import numpy as np
import wfdbx

wfdbx.write_record("pat", np.zeros(360, dtype=int), fs=360, gains=200, adcres=12,
                   units="mV", names="ECG", comments=["recorded at rest"],
                   metadata={"Age": "63", "Sex": "F"})
print(wfdbx.read_metadata("pat"))      # [('Age', '63'), ('Sex', 'F')]
print(wfdbx.read_header("pat").comments)
# ['recorded at rest', ' Age: 63', ' Sex: F']: the free comment is not an entry
\end{lstlisting}

\apientry{get\_metadata}{Return the value of a metadata entry}
\begin{lstlisting}[language=Python,style=signature]
get_metadata(record, key, *, input_dir=None, network=None)
\end{lstlisting}
Returns the value of one metadata entry of the header of a record.
\begin{params}
\item[record] Record name (as for \code{read\_metadata}).
\item[key] The key, such as \code{"Age"}. Letter case is ignored, and spaces at either end
  are removed.
\item[input\_dir] The input directory for this call only (optional; as for
  \code{read\_metadata}): it applies to relative record names only. Default \code{None}:
  the current setting.
\item[network] \code{True} to allow reading URLs for this call only, \code{False} to
  forbid it. Default \code{None}: the current setting.
\end{params}
\returns{The value, a string (\code{""} for an entry with an empty value), or
\code{None} if the header has no entry for \code{key}. If it has more than one, the first
is returned, and a \code{UserWarning} is issued.}
\raises{\code{IOError} if the header cannot be found or read; \code{ValueError} if
\code{record} is empty or \code{"-"}, or if \code{key} is not a valid key;
\code{TypeError} if \code{key} is not a string (or \code{bytes}).}
\begin{lstlisting}[language=Python,style=example]
import numpy as np
import wfdbx

wfdbx.write_record("pat", np.zeros(360, dtype=int), fs=360, gains=200, adcres=12,
                   units="mV", metadata={"Age": "63", "Sex": "F"})
print(wfdbx.get_metadata("pat", "age"))      # 63
print(wfdbx.get_metadata("pat", "Height"))   # None
\end{lstlisting}

\apientry{set\_metadata}{Set the value of a metadata entry}
\begin{lstlisting}[language=Python,style=signature]
set_metadata(record, key, value, *, input_dir=None, network=None, dir=None)
\end{lstlisting}
Sets the value of a metadata entry in the header of a record. If the header has an entry
for \code{key} (ignoring letter case), its line is replaced by \code{\# Key: value}, at
the same position, with the spelling of \code{key} given here; otherwise the entry is
added after the last info string of the header (or, if it has none, after its last line
that is not blank). The entry is written as \code{\# Key: value}, or \code{\# Key:} if
the value is empty. Every other line is kept, byte for byte, including its line ending.
\begin{params}
\item[record] Record name (as for \code{read\_metadata}). The header is changed where it
  was read, unless \code{dir} is given; a header read from a URL can only be written
  elsewhere, with \code{dir}.
\item[key] The key, such as \code{"Age"} or \code{"Date of birth"}. Spaces at either end
  are removed. It must not be empty, and must not contain \code{:}, \code{\#}, a line
  break or a NUL character.
\item[value] The value: a string, which may be a sentence. Spaces at either end are
  removed. It must not contain a line break or a NUL character; it may be empty.
\item[input\_dir] The input directory in which the header is read, for this call only
  (optional; as for \code{read\_metadata}): it applies to relative record names only.
  Default \code{None}: the current setting.
\item[network] \code{True} to allow reading the header from a URL, for this call only,
  \code{False} to forbid it. Default \code{None}: the current setting.
\item[dir] A directory, which must exist, in which to write the changed header, under the
  same file name, instead of changing the header itself; the signal files are not copied.
  Default \code{None}: the header itself is changed.
\end{params}
\returns{\code{None}.}
\raises{\code{IOError} if the header cannot be read, or the changed header cannot be
written; \code{ValueError} if \code{key} or \code{value} is not valid, if the header has
more than one entry for \code{key}, if it has no record line, if it was read from a URL
and \code{dir} is not given, or if \code{record} is empty or \code{"-"};
\code{TypeError} if \code{key} or \code{value} is not a string (or \code{bytes});
\code{FileNotFoundError} if \code{dir} does not exist.}
\begin{lstlisting}[language=Python,style=example]
import os
import numpy as np
import wfdbx

wfdbx.write_record("pat", np.zeros(360, dtype=int), fs=360, gains=200, adcres=12,
                   units="mV", metadata={"Age": "63", "Sex": "F"})
wfdbx.set_metadata("pat", "AGE", "64")              # replaces "# Age: 63"
wfdbx.set_metadata("pat", "Diagnosis", "sinus rhythm, no ectopy")   # a new entry
print(wfdbx.read_metadata("pat"))
# [('AGE', '64'), ('Sex', 'F'), ('Diagnosis', 'sinus rhythm, no ectopy')]

# Write the changed header to another directory; pat.hea is left unchanged.
os.makedirs("changed", exist_ok=True)
wfdbx.set_metadata("pat", "Sex", "M", dir="changed")
print(wfdbx.get_metadata("pat", "Sex"), wfdbx.get_metadata("changed/pat", "Sex"))   # F M
\end{lstlisting}

\apientry{remove\_metadata}{Remove a metadata entry}
\begin{lstlisting}[language=Python,style=signature]
remove_metadata(record, key, *, input_dir=None, network=None, dir=None)
\end{lstlisting}
Removes every entry for \code{key} (ignoring letter case) from the header of a record, so
it also repairs a header that has more than one. Other lines are kept, byte for byte.
Nothing is written if the header has no entry for \code{key}, unless \code{dir} is given.
\begin{params}
\item[record] Record name (as for \code{set\_metadata}).
\item[key] The key (as for \code{set\_metadata}).
\item[input\_dir] The input directory in which the header is read, for this call only
  (optional; as for \code{read\_metadata}): it applies to relative record names only.
  Default \code{None}: the current setting.
\item[network] \code{True} to allow reading the header from a URL, for this call only,
  \code{False} to forbid it. Default \code{None}: the current setting.
\item[dir] A directory, which must exist, in which to write the changed header instead of
  changing the header itself. Default \code{None}: the header itself is changed.
\end{params}
\returns{The number of lines removed (an \code{int}; 0 if the header has no entry for
\code{key}).}
\raises{As \code{set\_metadata}, except that more than one entry for \code{key} is not an
error.}
\begin{lstlisting}[language=Python,style=example]
import numpy as np
import wfdbx

wfdbx.write_record("pat", np.zeros(360, dtype=int), fs=360, gains=200, adcres=12,
                   units="mV", metadata={"Age": "63", "Sex": "F"})
print(wfdbx.remove_metadata("pat", "sex"))      # 1
print(wfdbx.remove_metadata("pat", "Sex"))      # 0: there is no such entry now
print(wfdbx.read_metadata("pat"))               # [('Age', '63')]
\end{lstlisting}

\apientry{update\_metadata}{Make several changes to the metadata}
\begin{lstlisting}[language=Python,style=signature]
update_metadata(record, changes, *, input_dir=None, network=None, dir=None)
\end{lstlisting}
Makes several changes to the metadata of a record, in one write of the header. Each
change is applied in order, as \code{set\_metadata} does, or, for a value of
\code{None}, as \code{remove\_metadata} does. All the changes are checked first: if one is
not valid, none is made.
\begin{params}
\item[record] Record name (as for \code{set\_metadata}).
\item[changes] A mapping from key to value, or a sequence of (key, value) pairs; a value
  of \code{None} removes the entry. No two keys may differ only in letter case.
\item[input\_dir] The input directory in which the header is read, for this call only
  (optional; as for \code{read\_metadata}): it applies to relative record names only.
  Default \code{None}: the current setting.
\item[network] \code{True} to allow reading the header from a URL, for this call only,
  \code{False} to forbid it. Default \code{None}: the current setting.
\item[dir] A directory, which must exist, in which to write the changed header instead of
  changing the header itself. Default \code{None}: the header itself is changed.
\end{params}
\returns{\code{None}.}
\raises{As \code{set\_metadata}; also \code{ValueError} if two keys of \code{changes}
differ only in letter case, and \code{TypeError} if \code{changes} is neither a mapping
nor a sequence of pairs.}
\begin{lstlisting}[language=Python,style=example]
import numpy as np
import wfdbx

wfdbx.write_record("pat", np.zeros(360, dtype=int), fs=360, gains=200, adcres=12,
                   units="mV", metadata={"Age": "63", "Sex": "F"})
wfdbx.update_metadata("pat", {"Age": "64", "Sex": None, "Dx": "164889003"})
print(wfdbx.read_metadata("pat"))      # [('Age', '64'), ('Dx', '164889003')]
try:
    wfdbx.update_metadata("pat", [("Dx", "426783006"), ("dx", None)])
except ValueError as e:
    print(e)        # metadata key 'dx' is given more than once (...)
print(wfdbx.get_metadata("pat", "Dx"))  # 164889003: nothing was changed
\end{lstlisting}

% ------------------------------------------------------------------------------------------
\subsection{The sample record and the example notebooks}
\label{sec:python-resource-functions}

\apientry{sample\_data\_dir}{The directory of the sample record}
\begin{lstlisting}[language=Python,style=signature]
sample_data_dir()
\end{lstlisting}
Returns the directory, within the installed package, that holds the sample record
\code{100s} (\cref{sec:python-resources}): its header, signal file and reference
annotations.
The function has no parameters.
\returns{The absolute path of the directory, a string.}
\raises{Nothing.}
\begin{lstlisting}[language=Python,style=example]
import os
import wfdbx

d = wfdbx.sample_data_dir()
print(sorted(f for f in os.listdir(d) if f.startswith("100s")))
# ['100s.atr', '100s.dat', '100s.hea']
rec = wfdbx.read_record("100s", sampto=10, input_dir=d)
print(rec.signals.shape)        # (10, 2)
\end{lstlisting}

\apientry{examples\_dir}{The directory of the example notebooks}
\begin{lstlisting}[language=Python,style=signature]
examples_dir()
\end{lstlisting}
Returns the directory, within the installed package, that holds the example notebooks
(\file{*.ipynb}, \cref{tab:python-notebooks}). The notebooks are meant to be copied
elsewhere (\code{copy\_examples}) before they are run.
The function has no parameters.
\returns{The absolute path of the directory, a string.}
\raises{Nothing.}
\begin{lstlisting}[language=Python,style=example]
import os
import wfdbx

print(sorted(os.listdir(wfdbx.examples_dir()))[0])     # 01_getting_started.ipynb
\end{lstlisting}

\apientry{copy\_examples}{Copy the example notebooks}
\begin{lstlisting}[language=Python,style=signature]
copy_examples(dest)
\end{lstlisting}
Copies the example notebooks to a directory, created if needed, where they can be opened
and run. Existing files of the same names are replaced. The command \cmd{wfdbx -{}-examples
DIR} does the same (\cref{sec:python-tools}).
\begin{params}
\item[dest] The directory into which the notebooks are copied (a path).
\end{params}
\returns{The list of the paths of the copied notebooks (\code{os.path.join(dest, name)}),
in the order of their names.}
\raises{\code{OSError} if the directory cannot be created or a notebook cannot be
written.}
\begin{lstlisting}[language=Python,style=example]
import os
import wfdbx

paths = wfdbx.copy_examples("notebooks")
print(len(paths), os.path.basename(paths[0]))    # 7 01_getting_started.ipynb
\end{lstlisting}

% ------------------------------------------------------------------------------------------
\subsection{The result classes}
\label{sec:python-results}

The readers return objects of the classes \code{Record}, \code{Header} and
\code{Annotations}, whose attributes are listed in
\cref{tab:python-record,tab:python-header,tab:python-annotations}; \code{check\_record}
returns a \code{CheckReport}, which holds \code{Issue} objects
(\cref{tab:python-checkreport,tab:python-issue}). The signals of a record are described
by \code{Siginfo} objects (\cref{tab:python-siginfo}), which are those of the WFDB library
(\cref{sec:python-library}). All these classes are available at the top level of the
package (\code{wfdbx.Record}, \ldots). The attributes of \code{Record}, \code{Header} and
\code{Annotations} are plain Python attributes, which may be changed.

\begin{xltabular}{\linewidth}{@{}lL@{}}
\caption{The attributes of \code{Record}, the result of \code{read\_record}}
\label{tab:python-record} \\
\toprule
\tabhead Attribute & \tabhead Contents \\
\midrule
\endfirsthead
\toprule
\tabhead Attribute & \tabhead Contents \\
\midrule
\endhead
\bottomrule
\endfoot
\code{name} & the record name, as given to \code{read\_record} \\
\code{fs} & the sampling frequency in Hz (a \code{float}), or 0 if the header gives none;
  with \code{highres=True}, the highest sampling frequency of a multi-frequency record \\
\code{signals} & a NumPy array with one row per sample and one column per signal read:
  \code{int32} ADC units, or \code{float64} physical units with \code{physical=True};
  invalid samples are \code{WFDB\_INVALID\_SAMPLE}, or NaN in physical units \\
\code{siginfo} & a list of \code{Siginfo} objects, one per signal read, in the order of
  the columns; \code{desc} has no white space at either end, and \code{adcres} is 0 if
  the header does not give it \\
\code{names} & the description of each signal, such as \code{"MLII"}, without the white
  space that a header may put at either end of it \\
\code{units} & the physical units of each signal, or \code{None} where the header gives
  none \\
\code{gains} & the gain of each signal, in ADC units per physical unit (\code{float}); 0
  means that the signal is uncalibrated \\
\code{baselines} & the ADC value of each signal that corresponds to 0 physical units \\
\code{base\_time} & the start time and date as text, to the second, such as
  \code{"[14:00:03 21/06/2024]"} (\code{"[14:00:03]"} without a date), or \code{None} if
  the header gives none \\
\code{base\_time\_ms} & the start time in milliseconds after midnight, such as
  \code{50403500} for 14:00:03.500, or \code{None} if \code{base\_time} is \code{None} \\
\code{base\_date} & the start date as a tuple \code{(year, month, day)}, or \code{None}
  if the header gives none \\
\code{sampfrom} & the number of the first sample returned \\
\code{comments} & the list of the info strings of the header (without the \code{\#}): its
  lines that begin with \code{\#} after the record line, found as the header parser finds
  it (\cref{sec:info-strings}) \\
\code{checksum\_ok} & \code{False} if the samples do not match the checksums of the
  header; checksums are checked only when a read reaches the end of the record, so a
  shorter read gives \code{True} \\
\end{xltabular}

\begin{xltabular}{\linewidth}{@{}lL@{}}
\caption{The attributes of \code{Header}, the result of \code{read\_header}}
\label{tab:python-header} \\
\toprule
\tabhead Attribute & \tabhead Contents \\
\midrule
\endfirsthead
\toprule
\tabhead Attribute & \tabhead Contents \\
\midrule
\endhead
\bottomrule
\endfoot
\code{name} & the record name, as given to \code{read\_header} \\
\code{fs} & the sampling frequency in Hz, or 0 if the header gives none \\
\code{counter\_freq} & the counter frequency, as \code{getcfreq()} returns it (the
  sampling frequency if the header gives none) \\
\code{base\_counter} & the base counter value, as \code{getbasecount()} returns it \\
\code{nsamples} & the number of samples per signal, or 0 if the header does not give it \\
\code{nsig} & the number of signals \\
\code{siginfo} & a list of \code{Siginfo} objects, one per signal; \code{desc} has no white
  space at either end, and \code{adcres} is 0 if the header does not give it \\
\code{base\_time}, \code{base\_time\_ms}, \code{base\_date} & the start time and date, as
  in \code{Record} \\
\code{comments} & the list of the info strings of the header (without the \code{\#}), as
  in \code{Record} \\
\code{metadata} & the list of the (key, value) pairs of the metadata entries of the
  header, in file order, as \code{read\_metadata} returns them \\
\end{xltabular}

\begin{xltabular}{\linewidth}{@{}lL@{}}
\caption{The attributes of \code{Annotations}, the result of \code{read\_annotations}.
  \code{len(ann)} is the number of annotations, which are in file order.}
\label{tab:python-annotations} \\
\toprule
\tabhead Attribute & \tabhead Contents \\
\midrule
\endfirsthead
\toprule
\tabhead Attribute & \tabhead Contents \\
\midrule
\endhead
\bottomrule
\endfoot
\code{record} & the record name, as given \\
\code{annotator} & the annotator name, such as \code{"atr"} \\
\code{samples} & the annotation times, in samples (an \code{int64} array) \\
\code{codes} & the annotation codes (an \code{int64} array), such as \code{wfdbx.NORMAL}
  or \code{wfdbx.PVC} \\
\code{symbols} & the list of the mnemonics of the codes, such as \code{"N"}, \code{"V"}
  or \code{"+"} \\
\code{subtypes} & the subtype of each annotation (an \code{int64} array) \\
\code{chans} & the channel number of each annotation (an \code{int64} array) \\
\code{nums} & the annotator number of each annotation (an \code{int64} array) \\
\code{aux} & the list of the auxiliary texts of the annotations (strings), with
  \code{None} where there is none \\
\code{fs} & the frequency in which \code{samples} are counted, normally the sampling
  frequency of the record; 0 if neither the record, nor the file, nor the argument
  \code{fs} of \code{read\_annotations} gives one \\
\end{xltabular}

\begin{table}[htbp]
\centering
\caption{The attributes and methods of \code{CheckReport}, the result of
  \code{check\_record}}
\label{tab:python-checkreport}
\begin{tabularx}{\linewidth}{@{}lL@{}}
\toprule
\tabhead Attribute & \tabhead Contents \\
\midrule
\code{record} & the record name, as given \\
\code{ok} & \code{True} if no issue has the severity \code{"error"} \\
\code{issues} & the list of the \code{Issue} objects, in the order of the checks \\
\code{lines()} & the text form of the report, as a list of lines: one line per issue
  (\code{str(issue)}), or the one line \code{no problems found} \\
\code{str(report)} & the lines of \code{lines()}, joined by line feeds (what
  \code{print(report)} prints) \\
\bottomrule
\end{tabularx}
\end{table}

\begin{table}[htbp]
\centering
\caption{The attributes of \code{Issue}, a problem found by \code{check\_record}}
\label{tab:python-issue}
\begin{tabularx}{\linewidth}{@{}lL@{}}
\toprule
\tabhead Attribute & \tabhead Contents \\
\midrule
\code{severity} & \code{"error"}, \code{"warning"} or \code{"info"} \\
\code{code} & a stable identifier of the kind of problem, such as \code{"checksum"} or
  \code{"gain\_missing"} (\cref{tab:python-issue-codes}) \\
\code{where} & where the problem is: \code{""} (the record), \code{"signal N"},
  \code{"segment SEG"}, \code{"segment SEG, signal N"} or \code{"annotator ANN"} \\
\code{message} & the description of the problem \\
\code{str(issue)} & \code{"SEVERITY: WHERE: MESSAGE"}, or \code{"SEVERITY: MESSAGE"}
  when \code{where} is empty \\
\bottomrule
\end{tabularx}
\end{table}
Two \code{Issue} objects are equal (\code{==}) when their four attributes are.

\begin{xltabular}{\linewidth}{@{}llL@{}}
\caption{The codes of the issues reported by \code{check\_record}}
\label{tab:python-issue-codes} \\
\toprule
\tabhead Code & \tabhead Severity & \tabhead Problem \\
\midrule
\endfirsthead
\toprule
\tabhead Code & \tabhead Severity & \tabhead Problem \\
\midrule
\endhead
\bottomrule
\endfoot
\multicolumn{3}{@{}l}{\textit{The header}} \\
\code{header\_unreadable} & error & the header cannot be read \\
\code{header\_line\_long} & error & a line of the header is longer than 255 characters \\
\code{header\_not\_ascii} & warning & a line of the header contains non-ASCII characters \\
\code{header\_name} & warning & the header names a record other than its file name \\
\code{fs\_missing} & warning & the sampling frequency is not specified \\
\code{length\_missing} & warning & the record length is not specified \\
\code{length\_large} & info & the record has more than 2147483647 samples \\
\code{no\_signals} & info & the record has no signals \\
\multicolumn{3}{@{}l}{\textit{The signal specifications}} \\
\code{desc\_missing} & warning & the signal has no description \\
\code{desc\_number} & warning & the description is only a number \\
\code{desc\_control} & error & the description contains control characters \\
\code{desc\_duplicate} & warning & several signals have the same description \\
\code{gain\_missing} & warning & the signal is uncalibrated (no gain is given) \\
\code{gain\_negative} & info & the gain is negative \\
\code{units\_missing} & warning & the units are not specified \\
\code{units\_control} & error & the units contain control characters \\
\code{adcres\_missing} & warning & the ADC resolution is not specified \\
\code{adcres\_too\_large} & error & the ADC resolution exceeds what the format can store \\
\code{adc\_range} & error & the ADC range does not fit in the format \\
\code{format\_difference} & warning & format 8 (first differences) cannot follow steep
  slopes \\
\code{filename} & warning & the signal file name contains characters other than letters,
  digits, \code{\_}, \code{-} and \code{.} \\
\multicolumn{3}{@{}l}{\textit{The signal files}} \\
\code{file\_unreadable} & error & a signal file cannot be read \\
\code{file\_large} & warning & a signal file has more than 2147483647 bytes \\
\code{read\_error} & error & the signal data cannot be read after some frames \\
\code{sample\_range} & error & samples are outside the ADC range \\
\code{invalid\_samples} & info & samples are invalid (missing data) \\
\code{checksum} & error & the checksum of a signal does not match its samples \\
\code{length\_short} & error & the signal files end before the length given in the
  header \\
\code{length\_long} & warning & the signal files hold more than the length given in the
  header \\
\code{length\_unspecified\_frames} & info & the number of frames that the signal files
  hold, when the header does not give the length \\
\multicolumn{3}{@{}l}{\textit{The segments of a multi-segment record}} \\
\code{segment\_unreadable} & error & a segment cannot be read \\
\code{nested} & error & a segment is itself a multi-segment record \\
\code{segment\_fs} & error & the sampling frequency of a segment differs from the
  record's \\
\code{segment\_lengths} & warning & the segment lengths do not add up to the length in
  the master header \\
\code{layout\_signals} & error & a segment has more signals than the record \\
\code{layout\_mismatch} & error & a signal of a segment differs from the record's
  signal \\
\code{layout\_missing\_signal} & error & a signal of a segment is not in the layout \\
\multicolumn{3}{@{}l}{\textit{The annotation files}} \\
\code{annotator\_name} & warning & the annotator name consists only of digits \\
\code{annotations\_unreadable} & error & the annotation file cannot be read \\
\code{annotations\_truncated} & error & the annotation file ends unexpectedly \\
\code{annotations\_order} & error & annotations are out of order \\
\code{annotation\_code} & error & annotations have an invalid code \\
\code{annotation\_null} & warning & annotations have code 0 (\code{NOTQRS}) \\
\code{annotation\_undefined} & warning & annotations have a code without a defined
  mnemonic \\
\code{annotations\_after\_end} & warning & annotations are after the end of the
  record \\
\end{xltabular}

\begin{xltabular}{\linewidth}{@{}lL@{}}
\caption{The attributes of \code{Siginfo}, the specification of a signal}
\label{tab:python-siginfo} \\
\toprule
\tabhead Attribute & \tabhead Contents \\
\midrule
\endfirsthead
\toprule
\tabhead Attribute & \tabhead Contents \\
\midrule
\endhead
\bottomrule
\endfoot
\code{fname} & the name of the signal file, as the header gives it \\
\code{desc} & the signal description, such as \code{"MLII"}, or \code{None} \\
\code{units} & the physical units, such as \code{"mV"}, or \code{None} if not specified \\
\code{gain} & the gain, in ADC units per physical unit (a \code{float}; 0: uncalibrated) \\
\code{initval} & the value of the first sample \\
\code{group} & the signal group number (the signals of one group share one file) \\
\code{fmt} & the storage format, such as \code{16}, \code{212} or \code{516} \\
\code{spf} & the number of samples per frame (more than 1 for the faster signals of a
  multi-frequency record) \\
\code{bsize} & the block size, in bytes (0 for ordinary files) \\
\code{adcres} & the ADC resolution, in bits (0: not specified) \\
\code{adczero} & the ADC output for an input of 0\,V \\
\code{baseline} & the ADC value that corresponds to 0 physical units \\
\code{nsamp} & the number of samples (0 if unknown) \\
\code{cksum} & the 16-bit checksum of the samples \\
\end{xltabular}

% ------------------------------------------------------------------------------------------
\subsection{Warnings and errors}
\label{sec:python-warnings}

\paragraph{Exceptions.} The high-level functions report errors by exceptions of the
standard classes (\cref{tab:python-exceptions}). Each function checks its arguments
before it creates or changes a file, so an invalid argument leaves the files as they
were; \code{write\_record} and \code{write\_annotations} also write through temporary
files, so that a failure while writing (a full disk, for example) leaves the existing
files as they were too. The library may also print a message on the standard error when
it cannot open a file (as in \code{init: can't open header for record REC}); \code{wfdbquiet()} turns these
messages off (\cref{sec:python-library}).

\begin{table}[htbp]
\centering
\caption{The exceptions of the high-level API}
\label{tab:python-exceptions}
\begin{tabularx}{\linewidth}{@{}lL@{}}
\toprule
\tabhead Exception & \tabhead Raised when \\
\midrule
\code{IOError} (\code{OSError}) & a header, signal file or annotation file cannot be
  found, read, created or written completely (for example on a full disk); a signal or
  annotation file ends before its header or its end-of-file marker says \\
\code{FileNotFoundError} & the directory given as \code{dir} does not exist (a subclass of
  \code{OSError}) \\
\code{ValueError} & an argument is invalid, or the operation needs a value that the
  record does not give (a time in seconds without a sampling frequency, physical units
  without a gain) \\
\code{TypeError} & an argument has the wrong type (a record or annotator name that is
  not a string in \code{check\_record}; a metadata key or value that is not a string;
  metadata that is neither a mapping nor a sequence of pairs) \\
\bottomrule
\end{tabularx}
\end{table}

\paragraph{The checksum warning.} When \code{read\_record} reads to the end of a record
whose samples do not match the checksums of its header, it returns all the samples and
issues a \code{RuntimeWarning}: ``checksum error in record REC: the signal data do not
match the checksums in its header''. \code{checksum\_warning=False} turns the warning off;
\code{checksum\_ok} of the result is \code{False} either way. A signal file that ends
early, or cannot be read, is an \code{IOError}, not a checksum warning, even when the
checksum is verified in its last, incomplete frame: no incomplete frame is returned.

\paragraph{The duplicate-key warning.} When a header holds two metadata entries whose keys
differ only in letter case (written by other software), \code{read\_metadata},
\code{get\_metadata} and \code{read\_header} use the first and issue a \code{UserWarning}:
``header of record REC has more than one entry for key KEY (lines L1, L2); the first is
used''. \code{set\_metadata} and \code{update\_metadata} then refuse to change that key
(\code{ValueError}), and \code{remove\_metadata} removes all of its entries.

Both warnings are issued with the module \code{warnings} of Python, so they can be
filtered, recorded or turned into errors (\code{warnings.simplefilter("error",
RuntimeWarning)}) as any other warning:

\begin{lstlisting}[language=Python,style=example]
import warnings
import numpy as np
import wfdbx

# A record whose last sample is changed after it was written: the checksum fails.
wfdbx.write_record("ck", np.arange(100), fs=100, gains=100, adcres=12, units="mV")
with open("ck.dat", "r+b") as f:
    f.seek(-2, 2)
    f.write(b"\x00\x00")
with warnings.catch_warnings(record=True) as caught:
    warnings.simplefilter("always")
    rec = wfdbx.read_record("ck")
print(rec.signals.shape, rec.checksum_ok)    # (100, 1) False: all the samples
print(caught[0].category.__name__)           # RuntimeWarning
rec = wfdbx.read_record("ck", checksum_warning=False)    # no warning

# Two entries for one key: the first is used, with a warning.
wfdbx.write_record("dup", np.zeros(10, dtype=int), fs=100, gains=0, adcres=12,
                   comments=[" Age: 63", " AGE: 64"])
with warnings.catch_warnings(record=True) as caught:
    warnings.simplefilter("always")
    print(wfdbx.get_metadata("dup", "age"))      # 63
print(caught[0].category.__name__)               # UserWarning
try:
    wfdbx.set_metadata("dup", "Age", "65")
except ValueError as e:
    print("ValueError:", e)
print(wfdbx.remove_metadata("dup", "age"))       # 2
\end{lstlisting}

\begin{lstlisting}[language=Python,style=example]
import numpy as np
import wfdbx

try:
    wfdbx.read_record("nosuchrecord")
except OSError as e:                  # IOError is OSError
    print("OSError:", e)              # OSError: can't read record nosuchrecord

wfdbx.write_record("raw", np.arange(10), fs=100, gains=0, adcres=12)   # uncalibrated
try:
    wfdbx.read_record("raw", physical=True)
except ValueError as e:
    print("ValueError:", e)           # ValueError: signal 0 of record raw is uncalibrated
\end{lstlisting}

% ==========================================================================================
\section{The WFDB library functions in Python}
\label{sec:python-library}

Every function of the WFDB library is available at the top level of the package under its
C name, with the same arguments, in the same order, and the same return values, so that
programs and documentation written for the C library carry over directly. The functions,
their arguments and their results are documented once, for all the languages, in
\cref{chap:library}; this section gives the rules by which the C declarations become
Python functions. \code{help(wfdbx.NAME)} gives the documentation of each function,
including what WFDB-x changes in it.

\subsection{Conventions}

\paragraph{Names.} The functions, the constants (\code{WFDB\_READ},
\code{WFDB\_INVALID\_SAMPLE}, \code{WFDB\_HIGHRES}, \ldots) and the annotation codes
(\code{NORMAL}, \code{PVC}, \ldots, \cref{chap:formats}) keep their C names. The macros of
\file{<wfdb/ecgmap.h>} are functions, under both their macro names and the names of the
functions behind them: \code{isqrs} and \code{wfdb\_isqrs}, \code{map1} and
\code{wfdb\_map1}, and so on. The structures are classes
(\cref{sec:python-structures}).

\paragraph{NULL is \code{None}.} A \code{NULL} pointer argument or result is
\code{None}: \code{sampfreq(None)} returns the frequency of the open record,
\code{getinfo(None)} continues with the info strings of the last record, and
\code{getinfo} returns \code{None} after the last one. Some trailing arguments have
defaults (\code{nann=None} and \code{nsig=None}: all the entries of the list):
\begin{lstlisting}[language=Python]
isigopen(record, siarray=None, nsig=0)
osigopen(record, siarray=None, nsig=0)
osigfopen(siarray, nsig=None)
setheader(record, siarray, nsig=None)
annopen(record, aiarray, nann=None)
\end{lstlisting}

\paragraph{Arrays and output arguments.} Where a C function fills an array that the caller
passes, the Python function fills a list that the caller passes:
\begin{itemize}
\item \code{isigopen} and \code{osigopen} empty the list \code{siarray} and append one
  \code{Siginfo} object (a copy) per signal; \code{siarray} may be \code{None}. With
  \code{nsig} $\le 0$, \code{isigopen} reads the header only, and returns the number of
  signals of the record; \code{osigopen(record, None, 0)} closes the open output signals,
  reads the header, opens nothing and returns 0.
\item \code{getvec}, \code{getframe} and \code{putvec} take a list (or a NumPy array) of
  integers, one per signal (\code{getframe} and \code{putvec}: one per sample of a
  frame); \code{getvec} and \code{getframe} replace its first entries, and it must have
  room for them (\code{[0] * nsig}), or \code{IndexError} is raised.
\item A C function that fills a structure through a pointer replaces the attributes of
  the object passed: \code{getann(n, annot)} fills an \code{Annotation},
  \code{getcal(desc, units, cal)} a \code{Calinfo}. \code{putann(n, annot)} and
  \code{putcal(cal)} read them.
\item An array of structures passed in is a list of objects: \code{annopen(record,
  [Anninfo("atr", WFDB\_READ)])}, \code{osigfopen([Siginfo(...), ...])}.
\item \code{getseginfo()} returns a list of \code{Seginfo} objects (empty if the open
  record is not a multi-segment record), instead of a pointer to an array.
\end{itemize}

\paragraph{Strings} are \code{str}. Text and binary data are handled as ISO-8859-1
(Latin-1), one character per byte, so every byte sequence round-trips unchanged. The
\code{aux} attribute of an annotation is \code{None}, or \code{bytes} whose first byte is
the length of the text that follows, as in C; \code{Annotation.aux\_string()} returns the
text (up to its first NUL byte), or \code{None}.

\paragraph{Return values and errors.} The functions return what the C functions return:
an \code{int} status (negative on errors), a number of signals, a \code{float}, a string
or \code{None}. Errors are reported by these values and by a message on the standard
error, as in C, never by exceptions (only a misuse of Python, such as a vector that is too
short, raises one). \code{wfdberror()} returns the last message; \code{wfdbquiet()} turns
the messages off and \code{wfdbverbose()} on again.

\paragraph{State.} As in C, the library keeps its state in the process (module-level
state): the open record, its input signals, the output signals and the annotators remain
open between calls, until \code{wfdbquit()} closes them. \code{wfdbquit()} leaves the
input and output directories and the network setting as they are, and so, as in C, the
time resolution of output annotators (\code{getafreq()}), which opening an EDF record sets
to its sampling frequency when none is set, and the getvec mode (\code{getgvmode()}), to
which it adds \code{WFDB\_HIGHRES}. The high-level functions of \cref{sec:python-api} call
\code{wfdbquit()}, which closes anything opened with the library functions, and leave
those two settings as they were.

\paragraph{Additions.} \code{getframes(n)} reads up to \code{n} frames at once, as
\code{n} calls of \code{getvec} would, but decoding whole blocks of the signal files. It
returns a tuple \code{(array, stat)}: an \code{int64} NumPy array with one row per frame
read and one column per signal, and the \code{getvec} status of the first frame that
could not be read, or the number of signals if all \code{n} frames were read.

\subsection{Module layout}

All the functions, classes and constants are imported by \code{import wfdbx}. They are
defined in these modules of the package:

\begin{table}[htbp]
\centering
\caption{The modules of the package \code{wfdbx}}
\label{tab:python-modules}
\begin{tabularx}{\linewidth}{@{}lL@{}}
\toprule
\tabhead Module & \tabhead Contents \\
\midrule
\code{wfdbx.\_wfdbio} & file lookup, the input and output directories, network access,
  error messages (\file{lib/wfdbio.c}); the settings functions \\
\code{wfdbx.\_signal} & headers, signal files in every format, time and date
  conversions, unit conversions (\file{lib/signal.c}); legacy mode \\
\code{wfdbx.\_annot} & annotation files, in the MIT and AHA formats
  (\file{lib/annot.c}) \\
\code{wfdbx.\_calib} & the calibration list (\file{lib/calib.c}) \\
\code{wfdbx.\_wfdbinit} & \code{wfdbinit}, \code{wfdbquit}, \code{wfdbflush}
  (\file{lib/wfdbinit.c}) \\
\code{wfdbx.\_ecgcodes} & the annotation codes and the mapping functions
  (\file{lib/ecgcodes.h}, \file{lib/ecgmap.h}) \\
\code{wfdbx.\_flac} & the FLAC codec of the formats 508, 516 and 524 \\
\code{wfdbx.\_cstd} & the exact behavior of the C standard library functions that the WFDB
  library relies on (\code{strtol}, \code{sscanf}, \code{printf}, \ldots) \\
\code{wfdbx.highlevel} & \code{read\_record}, \code{read\_header},
  \code{read\_annotations}, \code{write\_record}, \code{write\_annotations} and their
  result classes \\
\code{wfdbx.metadata} & the header metadata functions \\
\code{wfdbx.check} & \code{check\_record}, \code{CheckReport}, \code{Issue} \\
\code{wfdbx.resources} & \code{sample\_data\_dir}, \code{examples\_dir},
  \code{copy\_examples} \\
\code{wfdbx.cli} & the command \code{wfdbx} \\
\code{wfdbx.apps} & the command-line tools, one module per tool
  (\code{wfdbx.apps.rdsamp}, \ldots) \\
\bottomrule
\end{tabularx}
\end{table}

The modules whose names begin with an underscore are internal: their functions are used
through the top level of the package (\code{wfdbx.isigopen}, not
\code{wfdbx.\_signal.isigopen}).

\subsection{The structures}
\label{sec:python-structures}

The structures of the C library are classes whose attributes have the names of the
fields of the C structures (\cref{tab:python-structures}). Their constructors take the
attributes as arguments, in this order, all with defaults:
\begin{lstlisting}[language=Python]
Siginfo(fname=None, desc=None, units=None, gain=0.0, initval=0, group=0, fmt=0, spf=1,
        bsize=0, adcres=0, adczero=0, baseline=0, nsamp=0, cksum=0)
Seginfo(recname="", nsamp=0, samp0=0)
Anninfo(name=None, stat=WFDB_READ)
Annotation(time=0, anntyp=0, subtyp=0, chan=0, num=0, aux=None)
Calinfo(low=0.0, high=0.0, scale=0.0, sigtype=None, units=None, caltype=0)
\end{lstlisting}
\code{Siginfo}, \code{Annotation} and
\code{Calinfo} have a method \code{copy()}; \code{Annotation} also has
\code{assign(other)}, which copies the attributes of another annotation, and
\code{aux\_string()}.

\begin{xltabular}{\linewidth}{@{}lL@{}}
\caption{The classes of the C structures}
\label{tab:python-structures} \\
\toprule
\tabhead Class (C structure) & \tabhead Attributes \\
\midrule
\endfirsthead
\toprule
\tabhead Class (C structure) & \tabhead Attributes \\
\midrule
\endhead
\bottomrule
\endfoot
\code{Siginfo} (\code{WFDB\_Siginfo}) & \code{fname}, \code{desc}, \code{units},
  \code{gain}, \code{initval}, \code{group}, \code{fmt}, \code{spf}, \code{bsize},
  \code{adcres}, \code{adczero}, \code{baseline}, \code{nsamp}, \code{cksum}
  (\cref{tab:python-siginfo}) \\
\code{Seginfo} (\code{WFDB\_Seginfo}) & \code{recname} (the name of the segment),
  \code{nsamp} (its number of samples), \code{samp0} (the number of its first sample in
  the record) \\
\code{Anninfo} (\code{WFDB\_Anninfo}) & \code{name} (the annotator name, the suffix of the
  file), \code{stat} (\code{WFDB\_READ} or \code{WFDB\_WRITE} for the MIT format,
  \code{WFDB\_AHA\_READ} or \code{WFDB\_AHA\_WRITE} for the AHA format) \\
\code{Annotation} (\code{WFDB\_Annotation}) & \code{time} (in sample intervals),
  \code{anntyp} (the annotation code), \code{subtyp}, \code{chan}, \code{num},
  \code{aux} (\code{None}, or \code{bytes} whose first byte is the length of the text
  that follows) \\
\code{Calinfo} (\code{WFDB\_Calinfo}) & \code{low}, \code{high} (the levels of the
  calibration pulse, in physical units), \code{scale} (physical units per cm),
  \code{sigtype} (matched against the beginning of a signal description), \code{units},
  \code{caltype} (\code{WFDB\_AC\_COUPLED} or \code{WFDB\_DC\_COUPLED}, or-ed with
  \code{WFDB\_CAL\_SQUARE}, \code{WFDB\_CAL\_SINE} or \code{WFDB\_CAL\_SAWTOOTH}) \\
\end{xltabular}

\subsection{The settings functions}
\label{sec:python-settings}

The functions of \cref{tab:python-settings} are additions of WFDB-x to the library: they
set where files are read and written, network access and legacy mode, for the high-level
functions, the library functions and the tools alike (\cref{chap:files,chap:behavior}),
and give the path under which a program opens a file that it reads or creates itself, as
the tools do (\code{wfdbx\_input\_path}, \code{wfdbx\_output\_path};
\cref{sec:tool-files}). Each \code{set} function returns the previous setting, so that it
can be restored. The
C function \code{setwfdb(p)} sets the input directory too (\code{p} must name at most one
directory), and \code{getwfdb()} returns it (\code{"."} if none is set); in legacy mode
they set and return the WFDB path, as in C.

\begin{xltabular}{\linewidth}{@{}lL@{}}
\caption{The settings functions}
\label{tab:python-settings} \\
\toprule
\tabhead Function & \tabhead Effect and result \\
\midrule
\endfirsthead
\toprule
\tabhead Function & \tabhead Effect and result \\
\midrule
\endhead
\bottomrule
\endfoot
\code{wfdbx\_set\_input\_dir(d)} & sets the input directory: a local directory, a
  \code{\textasciitilde} path or a URL; \code{None} or \code{""} restores the default
  (\code{WFDBX\_INPUT\_DIR}, or else the current directory). Returns the previous input
  directory, or \code{None} if none was set \\
\code{wfdbx\_get\_input\_dir()} & returns the input directory, or \code{None} if none is
  set \\
\code{wfdbx\_set\_output\_dir(d)} & sets the output directory: a local directory or a
  \code{\textasciitilde} path (a URL is refused, with a message); \code{None} or
  \code{""} restores the default (\code{WFDBX\_OUTPUT\_DIR}, or else the current
  directory). Returns the previous output directory, or \code{None} (the current one if
  \code{d} is refused) \\
\code{wfdbx\_get\_output\_dir()} & returns the output directory, or \code{None} if none
  is set \\
\code{wfdbx\_set\_network(on)} & turns network access on (\code{True} or 1) or off
  (\code{False} or 0); \code{None} or a negative number restores the default (on if
  \code{WFDBX\_NETWORK} is a nonzero number). Returns the previous setting, 1 or 0 \\
\code{wfdbx\_get\_network()} & returns 1 if network access is on, 0 otherwise \\
\code{wfdbx\_set\_legacy(on)} & turns legacy mode on (\code{True} or 1) or off, for the
  whole library. Returns the previous setting, 1 or 0 \\
\code{wfdbx\_get\_legacy()} & returns 1 if legacy mode is on, 0 otherwise \\
\code{wfdbx\_legacy\_args(argv)} & if \code{argv[1]} (\code{argv}: the argument list,
  program name first) is \code{"-{}-legacy"}, removes it from the list, turns legacy mode on
  and returns 1; otherwise returns 0 \\
\code{wfdbx\_time\_needs\_freq(string)} & returns 1 if \code{strtim(string)} needs a
  sampling frequency, 0 otherwise \\
\code{wfdbx\_write\_errors()} & returns the number of write errors since the program
  started (output files that could not be written, flushed or closed) \\
\code{wfdbx\_note\_write\_error()} & counts a write error, for a program that writes a
  file itself \\
\code{wfdbx\_input\_path(name)} & returns the path under which a program opens a file
  that it reads itself: in normal mode, a relative name under the input directory, a
  \code{\textasciitilde} name expanded, an absolute name as given; in legacy mode, the
  name as given. \code{None}, \code{""} and \code{"-"} are returned as given; \code{None}
  for a URL, or a name that becomes one under the input directory \\
\code{wfdbx\_output\_path(name)} & returns the path under which a program creates a file
  that it writes itself: in normal mode, a relative name under the output directory (whose
  missing directories are then created), a \code{\textasciitilde} name expanded, an
  absolute name as given; in legacy mode, any name under the output directory, as a plain
  prefix. \code{None}, \code{""} and \code{"-"} are returned as given; \code{None} for a
  URL \\
\end{xltabular}

\subsection{Examples}

Reading the samples and the annotations of a record, as a C program does (the record
\code{100s} is read from the package's data directory):

\begin{lstlisting}[language=Python,style=example]
import wfdbx

previous = wfdbx.wfdbx_set_input_dir(wfdbx.sample_data_dir())

si = []                                  # receives one Siginfo per signal
nsig = wfdbx.isigopen("100s", si, 2)     # open 2 signals
print(nsig, [(s.desc, s.fmt, s.gain) for s in si])   # 2 [('MLII', 212, 200.0), ...]
print(wfdbx.sampfreq(None))              # 360.0: the frequency of the open record

v = [0] * nsig                           # receives one sample per signal
wfdbx.isigsettime(wfdbx.strtim("0:30"))  # skip to 30 seconds
wfdbx.getvec(v)
print(v, [wfdbx.aduphys(k, v[k]) for k in range(nsig)])

frames, stat = wfdbx.getframes(360)      # the next second, as a NumPy array
print(frames.shape, stat)                # (360, 2) 2

a = wfdbx.Annotation()
wfdbx.annopen("100s", [wfdbx.Anninfo("atr", wfdbx.WFDB_READ)])
n = 0
while wfdbx.getann(0, a) == 0:           # 0: an annotation was read; -1: the end
    if n < 2:
        print(wfdbx.mstimstr(a.time), wfdbx.annstr(a.anntyp), a.aux_string())
    n += 1
print(n, "annotations")                  # 75 annotations

wfdbx.wfdbquiet()                        # no messages on the standard error
print(wfdbx.isigopen("nosuchrecord", None, 0), wfdbx.wfdberror().strip())
wfdbx.wfdbverbose()

wfdbx.wfdbquit()                         # close everything
wfdbx.wfdbx_set_input_dir(previous)      # restore the input directory
\end{lstlisting}

Writing a record and an annotation file, sample by sample and annotation by annotation:

\begin{lstlisting}[language=Python,style=example]
import wfdbx

out = [wfdbx.Siginfo(fname="ramp.dat", desc="ramp", units="mV", gain=200.0, fmt=16,
                     adcres=12)]
print(wfdbx.osigfopen(out))         # 1: one output signal
wfdbx.setsampfreq(250.0)
for k in range(500):
    wfdbx.putvec([k])
print(wfdbx.newheader("ramp"))      # 0: ramp.hea is written
wfdbx.annopen("ramp", [wfdbx.Anninfo("qrs", wfdbx.WFDB_WRITE)])
for t in (100, 350):
    wfdbx.putann(0, wfdbx.Annotation(time=t, anntyp=wfdbx.NORMAL))
wfdbx.wfdbquit()                    # completes ramp.dat and ramp.qrs

print(open("ramp.hea").read().splitlines()[0])        # ramp 1 250 500
print(wfdbx.read_annotations("ramp", "qrs").samples)  # [100 350]
\end{lstlisting}

% ==========================================================================================
\section{The command-line tools in Python}
\label{sec:python-tools}

The 33 command-line tools (\code{rdsamp}, \code{rdann}, \code{wrsamp}, \code{xform},
\ldots) are run through the command \cmd{wfdbx}, installed with the package, followed by
the name of the tool and its arguments. The tools, their options, their output and their
exit status are those of the C tools, and are documented in \cref{chap:tools}.

\begin{lstlisting}[style=shell]
export WFDBX_INPUT_DIR=$(python -c "import wfdbx; print(wfdbx.sample_data_dir())")
wfdbx rdsamp -r 100s -t 10 -p     # the first 10 seconds of 100s, in physical units
wfdbx rdann -r 100s -a atr        # its annotations
wfdbx rdsamp --legacy -r 100s     # legacy mode: --legacy as the first argument
\end{lstlisting}

The command \cmd{wfdbx} also takes these options, instead of a tool:

\begin{table}[htbp]
\centering
\caption{The options of the command \code{wfdbx}}
\label{tab:python-wfdbx-options}
\begin{tabularx}{\linewidth}{@{}lL@{}}
\toprule
\tabhead Command & \tabhead Effect \\
\midrule
\cmd{wfdbx TOOL [ARGUMENTS ...]} & runs the tool \code{TOOL}, and exits with its
  status \\
\cmd{wfdbx -{}-list} & lists the names of the tools, one per line \\
\cmd{wfdbx -{}-examples [DIR]} & copies the example notebooks to the directory
  \code{DIR} (default \file{wfdbx-examples}), and prints their paths \\
\cmd{wfdbx -{}-version} & prints the version, \code{WFDB-x 1.2.2} \\
\cmd{wfdbx -{}-help}, \cmd{wfdbx -h} & prints the usage \\
\bottomrule
\end{tabularx}
\end{table}

Without arguments, \cmd{wfdbx} prints the usage and exits with status 1; an unknown tool
is an error (status 1). When an output file or the standard output cannot be written (a
full disk, for example), a tool reports it and exits with status 1.

\paragraph{Without the \cmd{wfdbx} command.} The package has no module
\code{wfdbx.\_\_main\_\_}, so \cmd{python -m wfdbx} is not supported. The command is the
module \code{wfdbx.cli}, and each tool is a module of \code{wfdbx.apps}, so that, with
\file{python/src} on \code{PYTHONPATH} for example:

\begin{lstlisting}[style=shell]
python -m wfdbx.cli rdsamp -r 100s -t 10     # as wfdbx rdsamp -r 100s -t 10
python -m wfdbx.apps.rdsamp -r 100s -t 10    # the tool's module itself
\end{lstlisting}

From Python, \code{wfdbx.cli.main(argv)} runs the command with the argument list
\code{argv} (without the program name), and the function \code{main(argv)} of a tool's
module runs the tool with \code{argv}, the complete argument list, program name first, as
in C. Both return the exit status; the tool writes to the standard output and error.

\begin{lstlisting}[language=Python,style=example]
import wfdbx.cli
from wfdbx.apps import rdsamp

status = wfdbx.cli.main(["rdsamp", "-r", "100s", "-t", "s2"])   # the first 2 samples
print(status)                                                   # 0
status = rdsamp.main(["rdsamp", "-r", "100s", "-f", "s10", "-t", "s11", "-p"])
\end{lstlisting}

\paragraph{Environment variables.} The tools, like the rest of the package, use the
environment variables of \cref{tab:python-environment}, also for the files that they read
or create themselves (\cref{sec:tool-files}): the \opt{-i} file of \cmd{wrsamp},
\cmd{edf2mit}, \cmd{wav2mit}, \cmd{a2m} and \cmd{ad2m}, the input files of
\cmd{ahaecg2mit} and the \opt{-S} script of \cmd{xform} are read from the input directory,
and the WAV file of \cmd{mit2wav}, the EDF file of \cmd{mit2edf}, the master header of
\cmd{snip -m} and the \file{.des}, \file{.ecg}, \file{.key} and \file{.txt} files of
\cmd{parsescp} are created in the output directory. Such a file cannot be read from a URL.
In legacy mode, the files that a tool reads itself are opened as named, and the output
directory, if set, is a plain prefix of every output name.

\begin{table}[htbp]
\centering
\caption{The environment variables used by the Python package}
\label{tab:python-environment}
\begin{tabularx}{\linewidth}{@{}lL@{}}
\toprule
\tabhead Variable & \tabhead Meaning \\
\midrule
\code{WFDBX\_INPUT\_DIR} & the input directory (a directory, a \code{\textasciitilde}
  path or a URL), when \code{wfdbx\_set\_input\_dir} does not set one \\
\code{WFDBX\_OUTPUT\_DIR} & the output directory, when \code{wfdbx\_set\_output\_dir} does
  not set one \\
\code{WFDBX\_NETWORK} & network access is on if it is a nonzero number, when
  \code{wfdbx\_set\_network} does not set it \\
\code{WFDBANNSORT} & \code{0} to keep annotation files in the order written (no sorting
  when they are closed) \\
\code{WFDBGVMODE} & the default getvec mode (\code{setgvmode} with a negative mode) \\
\code{WFDBCAL} & the calibration file read by \code{calopen(None)} (default
  \file{wfdbcal}) \\
\code{WFDB} & the WFDB path, used in legacy mode only (outside legacy mode, it is
  ignored, with a message) \\
\code{WFDB\_FLAC\_CHECK\_MD5} & if set, the MD5 signature of FLAC signal files is
  verified when they are read \\
\bottomrule
\end{tabularx}
\end{table}

% ==========================================================================================
\section{Testing}
\label{sec:python-testing}

The tests of the Python package are in \file{python/tests}. From the root of the
repository:

\begin{lstlisting}[style=shell]
pip install -e "./python[test]"       # the package and pytest, nbformat, nbclient, ipykernel
pytest python/tests
\end{lstlisting}

This runs the unit tests of the library, the high-level functions, the metadata functions
(with the cases shared by the seven languages, in \file{tests/metadata}) and
\code{check\_record}; it runs the example notebooks (when nbformat, nbclient and ipykernel
are installed); and it runs the cross-language conformance suite, which compares the
output of the tools with that of the original C programs, case by case (1155 cases, a few
minutes). The tests run against \file{python/src} when they are run from the repository,
in a temporary copy of the test records of \file{tests/data}. The tests that write to
\file{/dev/full} (write errors) are skipped where it does not exist, on macOS and Windows.

The environment variable \code{WFDBX\_CONFORMANCE} selects the conformance cases that
\cmd{pytest} runs: \code{0} skips the suite, and a list of tools, such as
\code{WFDBX\_CONFORMANCE=rdsamp,rdann}, runs the cases of those tools only.

\begin{lstlisting}[style=shell]
WFDBX_CONFORMANCE=0 pytest python/tests            # the unit tests only
WFDBX_CONFORMANCE=rdsamp,rdann pytest python/tests/test_conformance.py
\end{lstlisting}

The conformance suite can also be run on its own, with more options (the number of
parallel jobs, the legacy suite of 886 cases, which runs each tool with \opt{-{}-legacy},
verbose output, selected cases):

\begin{lstlisting}[style=shell]
PYTHONPATH=$PWD/python/src python3 tests/conformance/harness.py check --impl python -j4
PYTHONPATH=$PWD/python/src python3 tests/conformance/harness.py check --impl python --legacy
\end{lstlisting}
